% Notas ocultas al proyectar. El guion docente contiene los tiempos detallados.
% Cada etiqueta identifica una slide y permite mantener el guion sincronizado.

% S01
{
  \setbeamertemplate{navigation symbols}{}
  \setbeamertemplate{footline}{}
  \begin{frame}[plain,label=portada]
    \titlepage
  \end{frame}
}

% Nueva apertura. Los tiempos se administran en el guion docente.
\section{Motivación}

\begin{frame}[label=problematica]{El rol económico del Estado peruano}
  \textbf{¿Cómo estudiar el rol del Estado peruano en el desarrollo económico, reduciendo la influencia de nuestras afinidades políticas?}\par
  \vspace{0.5cm}
  Necesitamos criterios explícitos y evidencia que otras personas puedan revisar.\par
  \vspace{0.4cm}
  ¿Qué observarías para distinguir un Estado empresario, regulador o facilitador?
  \note{Recoger propuestas sin clasificar presidentes por reputación. Los procedimientos transparentes permiten discutir sesgos; no prometen eliminarlos. Esta pregunta amplia motiva un ejercicio acotado sobre declaraciones.}
\end{frame}

\begin{frame}[label=proxy-presidente]{Busquemos un indicador observable}
  \begin{columns}[T,onlytextwidth]
    \begin{column}{0.55\textwidth}
      Una \textbf{variable proxy} aproxima algo que no podemos observar directamente.\par
      \vspace{0.35cm}
      ¿Quién personifica a la nación?\par
      \vspace{0.25cm}
      El presidente de la República, según el artículo 110 de la Constitución vigente.
    \end{column}
    \begin{column}{0.40\textwidth}
      \centering
      \IfFileExists{images/presidente.png}{%
        \includegraphics[width=\linewidth,height=4.2cm,keepaspectratio]{presidente.png}%
      }{%
        \vspace{0.7cm}\textbf{Presidente de la República}\par
        \vspace{0.4cm}\textit{Espacio para tu ilustración}\par
        \vspace{0.3cm}{\small\texttt{images/presidente.png}}%
      }
    \end{column}
  \end{columns}
  \fuente{Constitución vigente, art. 110: \href{https://www3.congreso.gob.pe/Docs/files/constitucion/constitucion2019/files/basic-html/page67.html}{archivo del Congreso}.}
  \note{La norma vigente sirve para motivar la figura presidencial; no se aplica retroactivamente al marco constitucional de todos los mensajes históricos. La proxy que analizaremos es el contenido de sus declaraciones, no la persona del presidente ni su popularidad. Una imagen animada en GIF se exporta a PNG para el PDF.}
\end{frame}

\begin{frame}[label=orientacion-declarada]{Lo que un presidente declara en economía}
  Sus mensajes ofrecen indicios de la orientación económica que presenta y de las prioridades que anuncia.\par
  \vspace{0.45cm}
  \textbf{Esa orientación puede ayudar a formular hipótesis sobre la actuación del Estado.}\par
  \vspace{0.4cm}
  Para comprobarlas necesitamos contrastar el discurso con decisiones, normas, presupuestos y resultados.\par
  \vspace{0.25cm}
  ¿Qué podría impedir que lo anunciado se convierta en una política efectiva?
  \note{El mensaje es una comunicación pública situada, no acceso directo a pensamientos privados. Intervienen otras instituciones, restricciones y decisiones. Los discursos se usarán como proxy de orientación declarada, no como prueba de comportamiento estatal o efectos causales.}
\end{frame}

\begin{frame}[label=canales-presidente]{¿Dónde habla un presidente?}
  \begin{columns}[T,onlytextwidth]
    \begin{column}{0.23\textwidth}\centering
      {\color{peacockblue}\fontsize{34}{40}\selectfont\faMicrophone}\par
      \vspace{0.25cm}Conferencias de prensa
    \end{column}
    \begin{column}{0.23\textwidth}\centering
      {\fontsize{34}{40}\selectfont\bfseries X}\par
      \vspace{0.25cm}Publicaciones en X
    \end{column}
    \begin{column}{0.23\textwidth}\centering
      {\color{peacockblue}\fontsize{34}{40}\selectfont\faFacebookF}\par
      \vspace{0.25cm}Facebook
    \end{column}
    \begin{column}{0.23\textwidth}\centering
      {\color{UWMathPinkL4}\fontsize{34}{40}\selectfont\faFile}\par
      \vspace{0.25cm}Discursos presidenciales
    \end{column}
  \end{columns}
  \vspace{0.7cm}
  \textbf{Hoy trabajaremos con discursos presidenciales.}\par
  \vspace{0.25cm}
  Al terminar: ¿cómo adaptarías el ejercicio a publicaciones en X?
  \note{Iconos esquemáticos de Font Awesome y rótulo X, no fotografías de acontecimientos ni reproducciones de publicaciones reales. Cada canal tiene público, formato y contexto propios; no combinar sin definir una selección comparable.}
\end{frame}

\begin{frame}[label=api-concepto]{¿Cómo pediríamos datos a una plataforma?}
  Una \textbf{API} es una interfaz con reglas para que un programa solicite datos a un servicio y reciba una respuesta.\par
  \vspace{0.3cm}
  \centering
  \begin{tikzpicture}[font=\sffamily\small,>=stealth]
    \node[draw=UWBlack,minimum width=2.2cm,minimum height=1.1cm,align=center] (programa) at (0,0) {Tu programa\\Python};
    \node[draw=UWBlack,minimum width=1.7cm,minimum height=1.1cm,align=center] (api) at (4,0) {API};
    \node[draw=UWBlack,minimum width=2.2cm,minimum height=1.1cm,align=center] (servicio) at (8,0) {Servicio\\de la plataforma};
    \draw[->] ([yshift=0.25cm]programa.east) -- node[above]{Solicitud} ([yshift=0.25cm]api.west);
    \draw[->] ([yshift=0.25cm]api.east) -- node[above]{Consulta} ([yshift=0.25cm]servicio.west);
    \draw[<-] ([yshift=-0.25cm]programa.east) -- node[below]{Respuesta} ([yshift=-0.25cm]api.west);
    \draw[<-] ([yshift=-0.25cm]api.east) -- node[below]{Datos} ([yshift=-0.25cm]servicio.west);
  \end{tikzpicture}\par
  \raggedright\vspace{0.4cm}
  Ejemplo conceptual: solicitar publicaciones autorizadas con texto, fecha e identificador.\par
  \vspace{0.2cm}
  ¿Qué necesitarías definir antes de contarlas?
  \fuente{Extensión posterior: \href{https://docs.x.com/x-api/posts/lookup/quickstart}{documentación oficial de X}. Acceso sujeto a autenticación y condiciones del servicio.}
  \note{Diagrama conceptual editable. No es un endpoint real ni una llamada en vivo. La disponibilidad depende de permisos y condiciones de la plataforma; no se solicitan claves ni se instala software en esta sesión. API no equivale a descargar todo lo visible en una web.}
\end{frame}

\begin{frame}[label=objetivos]{Objetivos de aprendizaje}
  \textbf{Al finalizar esta clase deberías poder:}
  \begin{itemize}
    \item Cargar textos y reconocer qué documento representa cada fila.
    \item Modificar una instrucción y explicar su efecto sobre el resultado.
    \item Comparar menciones considerando la extensión de los documentos.
    \item Interpretar palabras en contexto y contrastar una etiqueta automática.
    \item Formular una conclusión con evidencia, límites y una aplicación a tu carrera.
  \end{itemize}
  \note{Presentar el producto común: tabla, dos fragmentos y una conclusión individual. API es motivación para una extensión posterior, no un objetivo de ejecución obligatorio.}
\end{frame}


% S02: 0--3
\begin{frame}[label=apertura]{Los discursos como fuente histórica}
\begin{columns}[T,onlytextwidth]
    \begin{column}{0.48\textwidth}
      {\bfseries ¿Cómo estudiar cientos de discursos?\par}
      \par\vspace{0.5cm}
      Propón un procedimiento.\par
      \par\vspace{0.3cm}
      ¿Qué contarías?\par
      ¿Qué comprobarías leyendo?
    \end{column}
    \begin{column}{0.45\textwidth}
      \centering
      \includegraphics[height=5.1cm]{28_07_1921_leguia_pagina1.png}\par
      \fuente{Mensaje de 1921, p. 1. Archivo del Congreso.}
    \end{column}
  \end{columns}
  \note{Recoger dos propuestas. Explicar que los mensajes permiten estudiar prioridades e instrumentos declarados. La actuación efectiva requiere otras fuentes.}
\end{frame}

% S03: 3--5. Extracto literal de u11860. Procedencia oculta hasta S08.
\begin{frame}[label=fragmento-a]{Fragmento A}
  \begin{quote}
    \enquote{El éxito alcanzado por la Empresa Estatal de Petróleo justifica por eso plenamente la decisión de autorizarle a emprender diversos e importantes programas, entre los que destacan los siguientes:}
  \end{quote}
  \par\vspace{0.3cm}
  ¿Qué intervención describe? Señala una expresión que sostenga tu interpretación.
  \note{Extracto literal de Velasco, 28 de julio de 1970, PDF p. 10, unidad u11860. Mostrar primero sin presidente. La unidad completa enumera ampliaciones de refinación y está en el notebook. No convertir la declaración de éxito en éxito verificado.}
\end{frame}

% S04: 5--7. Segunda oración literal de u14275.
\begin{frame}[label=fragmento-b]{Fragmento B}
  \begin{quote}
    \enquote{Esto supone que el Estado concentre su acción en el diseño de políticas generales y en la atención de servicios básicos y traslade al sector privado, con algunas excepciones, la responsabilidad del desarrollo de la producción y la comercialización de bienes y servicios, el que se dará bajo condición de competencia y de mercado.}
  \end{quote}
  ¿Qué cambia respecto de A? ¿Qué información falta?
  \note{Extracto literal de Alberto Fujimori, 28 de julio de 1991, PDF p. 10, unidad u14275. La oración anterior usa ente promotor. El notebook incluye ambas oraciones para discutir la ambigüedad. Mantener ocultas las etiquetas.}
\end{frame}

% S05: 7--10
\begin{frame}[label=pregunta]{Intervención económica en los discursos}
  \relax
{\bfseries ¿Cómo se describe la intervención económica del Estado?\par}
  \par\vspace{0.45cm}
  Compararemos palabras y leeremos los pasajes donde aparecen.\par
  \par\vspace{0.4cm}
  \textbf{Lo declarado en un discurso}\par
  Para comprobar lo que se ejecutó necesitamos otras fuentes.\par
  \par\vspace{0.4cm}
  Tu conclusión debe unir un resultado numérico y evidencia textual.
\end{frame}

% S06: 10--15
\begin{frame}[label=historia]{Enfoques de análisis textual}
\textbf{Análisis de contenido}\par
  Categorías y reglas para interpretar. Berelson, 1952.\par
  \par\vspace{0.3cm}
  \textbf{Conteos y diccionarios}\par
  Comparar palabras a escala. General Inquirer.\par
  \par\vspace{0.3cm}
  \textbf{Modelos de lenguaje}\par
  Clasificar mediante instrucciones. Gilardi y coautores, 2023.\par
  \par\vspace{0.35cm}
  ¿Qué aporta cada enfoque a nuestra pregunta?
  \fuente{\parencite{berelson1952,inquirer,gilardi2023}}
  \note{Cinco minutos como máximo. Los modelos de lenguaje pertenecen a los métodos computacionales. No presentar una sucesión que elimina lectura y categorías. La referencia bibliográfica de Berelson se corroboró con una reseña contemporánea, sin asumir lectura integral del libro. El estudio de Gilardi y coautores no valida el corpus de esta clase.}
\end{frame}

% S07: 15--18
\begin{frame}[label=roles]{Instrumentos económicos}
\textbf{Empresario}\par
  El Estado posee y controla la producción.\par
  \par\vspace{0.4cm}
  \textbf{Regulador}\par
  Establece reglas generales y supervisa mercados.\par
  \par\vspace{0.4cm}
  \textbf{Facilitador}\par
  Promueve actividades de forma selectiva, sin requerir propiedad estatal.\par
  \par\vspace{0.4cm}
  ¿Qué instrumento encuentras en A y B?
  \fuente{Categorías del proyecto. También distingue asistencial y no aplica.}
  \note{Son categorías del caso de investigación. Pedir que identifiquen el instrumento y su finalidad declarada. Evitar tratar cualquier ley como regulación económica o cualquier promoción como facilitación selectiva.}
\end{frame}

% S08: 18--23
\begin{frame}[label=decision-manual]{Clasificación y procedencia}
  Antes de revelar la procedencia, registra para A y B:
  \begin{enumerate}
    \item Una clasificación provisional.
    \item La expresión que la respalda.
    \item Una duda o información que falta.
  \end{enumerate}
  \par\vspace{0.3cm}
  Puedes justificar evidencia insuficiente. Conserva tu decisión para revisarla al final.
  \note{Volver a S03 y S04 o leer las unidades completas en el notebook. Revelar oralmente solo después de recoger respuestas: A, Velasco 1970, u11860, p. 10; B, Alberto Fujimori 1991, u14275, p. 10. Las etiquetas permanecen ocultas hasta S23.}
\end{frame}

\section{Práctica}

% S09: 23--26
\begin{frame}[label=entorno]{Python, bibliotecas y Colab}
  \renewcommand{\arraystretch}{1.4}
  \begin{tabular}{@{}p{0.25\textwidth}p{0.69\textwidth}@{}}
    \textbf{Python} & Lenguaje para escribir instrucciones.\\
    \textbf{Biblioteca} & Herramientas reutilizables para una necesidad.\\
    \textbf{Colab} & Entorno donde editamos el notebook y ejecutamos Python.\\
  \end{tabular}\par
  \par\vspace{0.4cm}
  Cada estudiante trabaja y guarda su propia copia. Puedes discutir con otras personas y ejecutar en tu laptop.
\end{frame}

% S10: 26--29
\begin{frame}[fragile,label=primera-celda]{Una variable y su resultado}
  Una variable conserva un valor. \texttt{=} asigna y las comillas delimitan texto.
  \begin{verbatim}
saludo = "Estoy explorando textos"
print(saludo)
  \end{verbatim}\par
  Cambia el saludo y vuelve a ejecutar. ¿Dónde aparece tu modificación?
  \note{El notebook usa un saludo algo más largo. La sintaxis es la misma. Mostrar también la edición de una celda de texto para guardar hipótesis.}
\end{frame}

% S11: 29--33
\begin{frame}[label=carga]{Carga de la tabla}
  \textbf{pandas} permite leer y explorar tablas.
  \par\vspace{0.3cm}
  \begin{enumerate}
    \item Ejecuta Preparación con la carga incluida.
    \item Ejecuta la tabla de documentos.
    \item Identifica fecha, orador y audiencia.
  \end{enumerate}
  ¿Qué representa una fila? ¿Todas tienen la misma extensión?
  \fuente{Documentación oficial de pandas: \parencite{pandas}.}
  \note{La carga incluida no necesita URL pública ni acceso a investigación. Si hace falta, indicar manual y subir datos estudiantes ZIP. La persona de apoyo atiende accesos mientras se explica la tabla.}
\end{frame}

% S12: 33--36
\begin{frame}[label=seleccion]{Seis mensajes seleccionados}
\begin{columns}[T,onlytextwidth]
    \begin{column}{0.31\textwidth}
      \centering
      \textbf{Leguía}\par
      1920 y 1921\par\vspace{0.15cm}
      \includegraphics[height=3.65cm]{28_07_1920_leguia_pagina1.png}
    \end{column}
    \begin{column}{0.31\textwidth}
      \centering
      \textbf{Velasco}\par
      1969 y 1970\par\vspace{0.15cm}
      \includegraphics[height=3.65cm]{28_07_1969_velasco_pagina1.png}
    \end{column}
    \begin{column}{0.31\textwidth}
      \centering
      \textbf{Alberto Fujimori}\par
      1991 y 1992\par\vspace{0.15cm}
      \includegraphics[height=3.65cm]{28_07_1991_alberto_fujimori_pagina1.png}
    \end{column}
  \end{columns}
  \par\vspace{0.2cm}
  Todos son mensajes del 28 de julio. Audiencias y ocasiones difieren.
  \fuente{Primera página de 1920, 1969 y 1991. Archivo oficial del Congreso; copias de \texttt{v1.0-ccl}. Regla de selección anterior al conteo.}
  \note{Seis documentos, 148 páginas y 619 unidades. Cada fila de esta tabla agrupada resume un par; cada fila del notebook representa un documento. No llamar homogénea ni representativa a la selección.}
\end{frame}

% S13: 36--43
\begin{frame}[label=transformacion]{Texto original y texto para conteos}
  \textbf{Original, extracto literal:}
  \begin{quote}
    \enquote{No ha descuidado mi Gobierno el incremento de la agricultura y ganadería.}
  \end{quote}
  \textbf{Versión para conteos:}
  \begin{quote}
    no ha descuidado mi gobierno el incremento de la agricultura y ganadería
  \end{quote}
  ¿Qué cambió? ¿Qué significa conservar \enquote{no}?
  \fuente{Leguía, 1921, PDF p. 11, unidad \texttt{u00126}.}
  \note{Mostrar la extracción y la versión preparada del notebook. La preparación además excluye rótulos iniciales, números de página y una nota editorial, y repara ciertos guiones. Definir token como cada aparición según la regla. La clase inspecciona ese procesamiento ya construido, sin ejecutar modelos lingüísticos.}
\end{frame}

% S14: 43--51
\begin{frame}[fragile,label=frecuencias]{Tabla de frecuencias}
  Antes de ejecutar, predice qué palabras aparecerán primero.
  \begin{verbatim}
numero_palabras = 10
  \end{verbatim}\par
  Cambia 10 a 8 y vuelve a ejecutar.
  \begin{itemize}
    \item \texttt{Counter} cuenta cada forma exacta.
    \item El ranking puede ocultar palabras frecuentes.
    \item El denominador conserva todas las palabras procesadas.
  \end{itemize}
  ¿Qué cambió en la tabla y qué permaneció igual?
\end{frame}

% S15: 51--54
\begin{frame}[label=prediccion]{Predicción sobre competencia}
  Compararemos Leguía de 1921 y Alberto Fujimori de 1991.
  \par\vspace{0.4cm}
  \begin{enumerate}
    \item ¿Cuál tendrá más menciones absolutas?
    \item ¿Cuánto puede explicar la extensión del mensaje?
    \item ¿La palabra describirá el mismo concepto?
  \end{enumerate}
  Escribe tu predicción antes de ejecutar. Puedes cambiar la palabra y los documentos.
\end{frame}

% S16: 54--58
\begin{frame}[label=normalizacion]{Menciones por 1.000 palabras}
  \[
    \text{Tasa} = 1\,000 \times
    \frac{\text{menciones de la palabra}}{\text{palabras del documento procesado}}
  \]
  \par\vspace{0.3cm}
  El denominador incluye las palabras ocultas solo del ranking.\par
  \par\vspace{0.3cm}
  La tasa describe presencia léxica. La clasificación de roles requiere interpretar las unidades.
  \note{Leer mil, no uno con tres decimales. Ambos conteos deben usar la misma regla. La tabla del notebook permite comprobar numerador y denominador. Una tasa de palabras no es porcentaje de unidades productivas.}
\end{frame}

% S17: 58--66. Mostrar solo después de las predicciones y la ejecución.
\begin{frame}[label=tasas]{Conteos y tasas de competencia}
\renewcommand{\arraystretch}{1.25}
  \begin{tabular}{@{}lrrr@{}}
    \textbf{Mensaje} & \textbf{Menciones} & \textbf{Palabras} & \textbf{Por 1.000}\\
    Leguía 1921 & 2 & 4.527 & 0,442\\
    Fujimori 1991 & 11 & 23.348 & 0,471\\
  \end{tabular}\par
  \par\vspace{0.12cm}
  \centering
  \begin{tikzpicture}
    \begin{axis}[
      xbar,xmin=0,xmax=0.6,width=0.85\textwidth,height=3.25cm,
      symbolic y coords={Leguía 1921,Fujimori 1991},ytick=data,
      enlarge y limits=0.7,bar width=11pt,axis lines*=left,
      xmajorgrids=true,grid style={draw=UWGrayL1},
      tick label style={font=\footnotesize},label style={font=\footnotesize},
      xticklabel style={/pgf/number format/use comma},
      xlabel={Menciones por 1.000 palabras}]
      \addplot[fill=UWGoldL3,draw=UWBlack] coordinates {
        (0.4417936823503424,{Leguía 1921})
        (0.4711324310433442,{Fujimori 1991})
      };
    \end{axis}
  \end{tikzpicture}\par
  \raggedright
  ¿Qué cambió al considerar la extensión?
  \fuente{Forma exacta: competencia. Datos didácticos; tasas redondeadas.}
  \note{La diferencia de conteos es mucho mayor que la de tasas. No declarar significación estadística ni inferir prioridad económica. El gráfico proyectado y el notebook usan los mismos valores; mostrar después de las predicciones y la ejecución. Si cambió la palabra, interpretar su resultado y usar esta slide únicamente como ejemplo de la consulta inicial. El gráfico de reserva está en el apéndice.}
\end{frame}

% S18: 66--74
\begin{frame}[label=recuperacion]{Dudas y recuperación}
\textbf{Si perdiste una variable}\par
  Ejecuta Preparación y después Recuperación de la comparación.\par
  \par\vspace{0.5cm}
  \textbf{Si reiniciaste la sesión}\par
  Sigue el mismo orden: Preparación y Recuperación.\par
  \par\vspace{0.5cm}
  Antes de continuar, comprueba tu tabla y su denominador.\par
  \par\vspace{0.2cm}
  Si el gráfico tarda, conserva la tabla para discutir el resultado.
  \note{Reservar ocho minutos incluidos en la agenda. La pausa breve cabe dentro de este bloque si las dificultades están resueltas. Una tabla guardada permite discutir, pero no sustituye ejecutar y modificar código en la laptop propia.}
\end{frame}

\section{Interpretación}

% S19: 74--77
\begin{frame}[label=contexto]{Palabras en contexto}
  ¿La palabra significa lo que supusimos?
  \par\vspace{0.4cm}
  Buscaremos en el texto original y conservaremos negaciones y signos.
  \begin{enumerate}
    \item Predice un uso de la palabra.
    \item Localiza sus coincidencias.
    \item Amplía la ventana y lee la unidad completa.
  \end{enumerate}
  \note{La ventana visible reemplaza saltos de línea por espacios. El contenido se conserva. La búsqueda literal puede diferir del conteo por los guiones reparados y rótulos excluidos. No interpretar cero coincidencias como ausencia del concepto.}
\end{frame}

% S20: 77--81
\begin{frame}[fragile,label=amplitud]{Amplitud del contexto}
  \begin{verbatim}
amplitud_contexto = 100
  \end{verbatim}\par
  Cambia 100 a 180. El número indica caracteres a cada lado.
  \par\vspace{0.3cm}
  Selecciona dos filas de documentos distintos y recupera sus unidades completas.
  \par\vspace{0.3cm}
  Guarda el identificador, la página y el fragmento. El visor del PDF permite comprobar el original.
  \note{Para la consulta inicial, filas 0 y 3. Si cambió la palabra o selección, los índices deben revisarse. El código deja vacía la evidencia cuando la selección es inválida para evitar conservar fragmentos antiguos.}
\end{frame}

% S21: 81--86. Mostrar después de que los estudiantes examinen sus pasajes.
\begin{frame}[label=polisemia]{Dos usos de competencia}
\begin{columns}[T,onlytextwidth]
    \begin{column}{0.47\textwidth}
      \textbf{Leguía, 1921}\par
      \begin{quote}\raggedright
        \enquote{cuya competencia se ha evidenciado en la campaña que, para extirpar la fiebre amarilla, desarrollóse en el norte de la República}
      \end{quote}
      \fuente{\texttt{u00126}, p. 11.}
    \end{column}
    \begin{column}{0.47\textwidth}
      \textbf{Alberto Fujimori, 1991}\par
      \begin{quote}\raggedright
        \enquote{el que se dará bajo condición de competencia y de mercado.}
      \end{quote}
      \fuente{\texttt{u14275}, p. 10.}
    \end{column}
  \end{columns}
  \par\vspace{0.5cm}
  ¿Qué significa competencia en cada pasaje?\par
  ¿Qué cambia en tu interpretación de los conteos?
  \note{Los dos fragmentos son ejemplos dirigidos, no una clasificación de todas las ocurrencias. Leer las unidades completas del notebook. La primera continúa con fomento agrícola, por lo que el segmento reúne distintos asuntos.}
\end{frame}

% S22: 86--90
\begin{frame}[label=conclusion-provisional]{Conclusión y evidencia}
  \begin{enumerate}
    \item ¿Qué diferencia numérica encontraste?
    \item ¿Qué muestran tus dos fragmentos?
    \item ¿Qué parte de tu hipótesis mantienes o revisas?
    \item ¿Qué afirmación restringe tu limitación?
  \end{enumerate}
  \par\vspace{0.3cm}
  Antes del contraste: resultado numérico y dos fragmentos identificados.
  \note{Aceptar resultados prudentes y los que no respaldan la primera hipótesis. Exigir el vínculo entre números y evidencia, sin atribuir un rol a todo un gobierno a partir de una palabra.}
\end{frame}

% S23: 90--95
\begin{frame}[label=etiquetas]{Contraste con etiquetas existentes}
  \renewcommand{\arraystretch}{1.4}
  \begin{tabular}{@{}lll@{}}
    \textbf{Ejemplo} & \textbf{Primaria del modelo} & \textbf{Confianza}\\
    A, Velasco 1970 & empresario & alta\\
    B, Fujimori 1991 & regulador & media\\
  \end{tabular}\par
  \par\vspace{0.4cm}
  ¿Coincides con el modelo?\par
  Justifica tu decisión con el instrumento descrito y el contexto.
  \fuente{Etiquetas existentes, \texttt{v1.0-ccl}. Confianza categórica, no probabilidad calibrada.}
  \note{Mostrar únicamente después de la clasificación manual. La secundaria de B es asistencial. Leer la justificación del modelo y conservar una duda si corresponde. Coincidir con él no demuestra exactitud. No ejecutar clasificación en vivo.}
\end{frame}

% S24: 95--98
\begin{frame}[label=limites]{Límites de las etiquetas}
  \begin{itemize}
    \item El segundo anotador fue otro modelo de lenguaje.
    \item La referencia de 30 unidades contiene un caso empresario y uno regulador.
    \item Las métricas por clase rara tienen poco soporte.
  \end{itemize}
  \par\vspace{0.3cm}
  La interpretación requiere revisión. Conteos, roles y actuación efectiva responden a preguntas diferentes.
  \fuente{Documentación de anotación y decisiones del proyecto, versión congelada.}
\end{frame}

\section{Cierre}

% S25: 98--106
\begin{frame}[label=entrega]{Tu análisis en el notebook}
\textbf{Resultado}\par
  Una tabla o gráfico con denominador y una conclusión.\par
  \par\vspace{0.3cm}
  \textbf{Evidencia}\par
  Dos fragmentos originales con documento y página.\par
  \par\vspace{0.3cm}
  \textbf{Alcance}\par
  Una limitación y su consecuencia; una aplicación a tu carrera.\par
  \par\vspace{0.35cm}
  Conserva tu predicción, el cambio de código y el contraste manual.\par
  \fuente{Entrega individual, sin calificación.}
\end{frame}

% S26: 102--106. Alternar con S25 según el avance.
\begin{frame}[label=aplicaciones]{Documentos de tu carrera}
  \renewcommand{\arraystretch}{1.3}
  \begin{tabular}{@{}p{0.32\textwidth}p{0.62\textwidth}@{}}
    \textbf{Área} & \textbf{Documentos y preguntas posibles}\\
    Derecho & Contratos y normas: obligaciones y excepciones.\\
    Economía & Noticias y discursos: instrumentos declarados.\\
    Finanzas & Reportes: cómo se describen riesgos.\\
    Administración e Ingeniería Empresarial & Reportes y reclamos: prioridades y problemas.\\
    Ingeniería de la Información & Encuestas abiertas: patrones y contexto.\\
  \end{tabular}\par
  \par\vspace{0.3cm}
  ¿Qué metadatos necesitarías para comparar tus documentos?
\end{frame}

% S27: 106--110
\begin{frame}[label=cierre]{Guardado y cierre}
\begin{enumerate}
    \item Guarda tu copia con respuestas y salidas visibles.
    \item Descarga el notebook.
    \item Súbelo a la carpeta de Drive indicada y comprueba la recepción.
  \end{enumerate}
  \par\vspace{0.3cm}
  {\small\texttt{apellido\_nombre\_analisis\_textual.ipynb}\par}
  \par\vspace{0.5cm}
  ¿Qué afirmación todavía no permiten tus datos?
  \note{106--109: guardar, descargar y subir. 109--110: cierre oral. La persona de apoyo registra incidencias. El enlace de entrega lo proporciona el docente; no inventar una URL.}
\end{frame}

\appendix
\section*{Apéndice}

% A01: consulta, fuera del recorrido común
\begin{frame}[label=denominadores]{Denominadores para preguntas distintas}
  \renewcommand{\arraystretch}{1.35}
  \begin{tabular}{@{}p{0.39\textwidth}p{0.55\textwidth}@{}}
    \textbf{Medida} & \textbf{Denominador}\\
    Tasa de palabra & Palabras procesadas del documento.\\
    Incidencia productiva & Todas las unidades oficiales seleccionadas.\\
    Mezcla de roles productivos & Solo unidades empresario, regulador y facilitador seleccionadas.\\
  \end{tabular}\par
  \par\vspace{0.3cm}
  Proyecto completo: 22.365 unidades oficiales y 6.899 productivas. Selección didáctica: 619 unidades. Sus resultados tienen alcances distintos.
  \note{No usar los seis ejemplos manuales dirigidos para calcular proporciones. Las etiquetas secundarias no duplican unidades. La entrega común no requiere calcular mezclas de roles.}
\end{frame}

% A02: gráfico editable nativo; valores completos del CSV incluido en el ZIP
\begin{frame}[label=grafico-reserva]{Gráfico de la consulta inicial}
  \centering
  \begin{tikzpicture}
    \begin{axis}[
      xbar, xmin=0, xmax=0.6, width=0.87\textwidth, height=4.8cm,
      symbolic y coords={Leguía 1921,Fujimori 1991}, ytick=data,
      enlarge y limits=0.7, bar width=13pt,
      axis lines*=left, xmajorgrids=true,
      grid style={draw=UWGrayL1},
      tick label style={font=\small}, label style={font=\small},
      xticklabel style={/pgf/number format/use comma},
      xlabel={Menciones por 1.000 palabras}
    ]
      \addplot[fill=UWGoldL3,draw=UWBlack] coordinates {
        (0.4417936823503424,{Leguía 1921})
        (0.4711324310433442,{Fujimori 1991})
      };
    \end{axis}
  \end{tikzpicture}
  \fuente{Forma exacta: competencia. Datos didácticos, consulta inicial.}
  \note{Usar como reserva después de las predicciones. El gráfico del notebook utiliza los mismos valores y un acento borgoña. Aquí se usa el amarillo del template. Gráfico editable en LaTeX, sin imagen rasterizada.}
\end{frame}

% A03
\begin{frame}[label=desafio]{Desafío opcional}
  Después de completar la entrega:
  \begin{enumerate}
    \item Predice qué pasará al retirar un documento.
    \item Ejecuta el desafío del notebook.
    \item Explica el cambio de selección y denominador.
  \end{enumerate}
  \par\vspace{0.3cm}
  La tasa conjunta divide la suma de menciones por la suma de palabras. No promedia tasas sin atender a la extensión.
\end{frame}

% A04
\begin{frame}[label=procedencia]{Procedencia de ejemplos y adaptaciones}
  \begin{itemize}
    \item A: Velasco 1970, \texttt{u11860}, PDF p. 10.
    \item B: Fujimori 1991, \texttt{u14275}, PDF p. 10.
    \item Polisemia: Leguía 1921, \texttt{u00126}, PDF p. 11.
  \end{itemize}
  \par\vspace{0.3cm}
  PDF públicos del Museo del Congreso. Copias, unidades y etiquetas del proyecto \texttt{v1.0-ccl}.
  \par\vspace{0.2cm}
  La progresión pedagógica adapta los materiales locales de análisis textual en R a una práctica propia en Python.
\end{frame}

% R01--R03: lectura posterior, sin tiempo adicional dentro de los 110 minutos
\begin{frame}[label=referencias-fundamentos]{Referencias del análisis textual}
  \small
  \nocite{grimmer2013}
  \printbibliography[heading=none,keyword=fundamentos]
\end{frame}

\begin{frame}[label=referencias-automatizacion]{Referencias de automatización}
  \small
  \printbibliography[heading=none,keyword=automatizacion]
\end{frame}

\begin{frame}[label=referencias-herramientas]{Documentación de herramientas}
  \small
  \nocite{matplotlib,colab}
  \printbibliography[heading=none,keyword=herramientas]
\end{frame}
